% !TEX root = approx8.tex

\subsection{Systems of categories with increasing
  accuracies} \label{s:sys-tpc}

In what follows we will deal with families of TPC's that should be
thought of as approximations of a (currently not yet defined) limit TPC. The
TPC's occurring in the family are parametrized by a space controlling
the accuracy-level of their approximation. This situation naturally
occurs when dealing with Fukaya categories that can be parametrized by
different choices of perturbation data. The setting and definitions
below were conceived with the example of Fukaya categories in
mind. Nevertheless we believe that a general axiomatic framework may
prove useful in other cases too.

We will axiomatize these notions in three different settings below,
starting with the most general one - a system of PC's.

\subsubsection{The case of PC's} \label{sb:sys-pc}
Let $(\mathcal{P}, \preceq)$ be a directed set (i.e.~a preordered set
such that
$\forall \; p', p'' \in \mathcal{P}, \; \exists \; q \in \mathcal{P}$
with $p', p'' \preceq q$), together with a function
$\nu: \mathcal{P} \longrightarrow \mathbb{R}_{\geq 0}$ that has the
following properties:
\begin{enumerate}
\item For every $p \preceq q$ we have $\nu(p) \geq \nu(q)$ (i.e.~$\nu$
  is decreasing with respect to $\preceq$).
\item For every $\delta>0$ there exists $p \in \mathcal{P}$ such that
  $\nu(p) \leq \delta$.
\end{enumerate}

We will refer to $(\mathcal{P}, \preceq , \nu)$ as the parameter space
for our system of PC's. For brevity we will omit $\preceq$ and $\nu$
from the notation and simply write $\mathcal{P}$ for the triple. We
will refer to $\nu(p)$ as the norm or size of the parameter
$p \in \mathcal{P}$.

Consider now a family
$\widehat{\msc} = \{\msc_p\}_{p \in \mathcal{P}}$ of PC's $\msc_p$
parametrized by $p \in \mathcal{P}$. We will denote the shift and
translation functors of all the $\msc_p$'s by the same notation,
$\Sigma$ and $T$, respectively. Assume further that for every
$p \preceq q$ we are given two sets
$\msj_{p,q} = \mathscr{J}_{p,q}(\widehat{\msc})$,
$\msj_{q,p} = \mathscr{J}_{q,p}(\widehat{\msc})$, of PC functors
$\msc_p \longrightarrow \msc_q$, $\msc_q \longrightarrow \msc_p$,
respectively, with the following properties: \label{pp:sys-functors}
\begin{enumerate}
\item $\msj_{p,p} = \{\id_{\msc_p}\}$ for every $p \in \mathcal{P}$.
\item \label{i:0-iso-sys-1} If $p \preceq q$ then all the functors in
  $\msj_{p,q}$ are mutually $0$-isomorphic. More specifically, for
  every two functors $\msf, \msg \in \msj_{p,q}$ there is a natural
  isomorphism $\msf \cong \msg$ of persistence level $0$ (in other
  words, $\msf$ and $\msg$ are isomorphic in the $0$-level category
  $\text{pc-fun}^0(\msc_p, \msc_q)$ of the PC of PC-functors
  $\text{pc-fun}(\msc_p, \msc_q)$). Here and in what follows we will
  abbreviate this by writing $\msf \cong_0 \msg$. Furthermore, we
  require the same thing also for any two functors in $\msj_{q,p}$.
\item \label{i:0-iso-sys-2} If $p \preceq q \preceq r$ then for every
  $\msf \in \msj_{p,q}$, $\msg \in \msj_{q,r}$, $\msh \in \msj_{p,r}$
  we have $\msg \circ \msf \cong_0 \msh$.
\item \label{i:d-int-func} If $p \preceq q$ then for every
  $\msf \in \msj_{p,q}$ and $\msh \in \msj_{q,p}$ we have
  $$\dint(\msh \circ \msf, \id_{\msc_p}), \, \dint(\msf \circ
  \msh, \id_{\msc_q}) \leq C | \nu(p) - \nu(q)|,$$ for some constant
  $C$ that depends only on the entire family $\widehat{\msc}$ but
  neither on $p, q$ nor on $\msf, \msg$. Here the two $\dint$'s
  stand for the interleaving distances on the PC's of PC functors
  $\msc_p \longrightarrow \msc_p$ and $\msc_q \longrightarrow \msc_q$.
\end{enumerate}
In what follows we will denote by $\widehat{\msc}$ the entire data
above, namely both the family of categories
$\{\msc_p\}_{p \in \mathcal{P}}$ as well the sets of functors
$\msj_{p,q}$ and refer to $\widehat{\msc}$ as a {\em system of PC's
  with increasing accuracy}, or sometimes a {\em system of PC's} for
short. We will call the functors in $\msj_{p,q}$ comparison functors
and typically denote such a functor by
$\msh_{p,q}: \msc_p \longrightarrow \msc_q$.


\begin{rem} \label{rem:pc-limits}
  \begin{enumerate}
  \item Given a system $\widehat{\msc}$ of PC's with increasing accuracy,
    all the $\infty$-level categories $\msc^{\infty}_p$, corresponding
    to the PC's $\msc_p$ in the system $\widehat{\msc}$, are mutually
    equivalent. From the persistence viewpoint, as $p$ ``grows'' (with
    respect to $\preceq$) the comparison between $\msc_p$ and
    $\msc_q$ for $p \preceq q$ becomes more and more
    precise. \label{i:rem-PC-1}
  \item A single PC $\msc$ makes a special case of the definition
    above. We can view it as a system of PC's parametrized by a
    $1$-point parameter space $\mathcal{P} = \{*\}$. In this case we
    set the norm function $\nu$ to be trivial, $\nu(*) =0$.
  \end{enumerate}
\end{rem}



\subsubsection{Functors of systems of PC's} \label{sbsb:func-sys-pc}
Next we discuss functors between systems of PC's with increasing
accuracy. Let $\widehat{\msc}$ and $\widehat{\msd}$ be two systems of
PC's with increasing accuracy. Whenever we need to distinguish between
the additional structures of each of the systems $\widehat{\msc}$ and
$\widehat{\msd}$ we will use $\widehat{\msc}$ and $\widehat{\msd}$
subscripts or superscripts, depending on notational convenience (for
example, the parameter space of $\widehat{\msc}$ will be denoted
$\mathcal{P}_{\widehat{\msc}}$, the comparison functors
$\msh_{p,q}^{\widehat{\msc}}$, etc. but we will denote the preorders
and size of parameters in both systems by $\preceq$ and $\nu$
respectively).

A functor
$\widehat{\msf} : \widehat{\msc} \longrightarrow \widehat{\msd}$
consists of the following structures:
\begin{enumerate}
\item A map
  $\widehat{\msf}: \mathcal{P}_{\widehat{\msc}} \longrightarrow
  \mathcal{P}_{\widehat{\msd}}$ between the parameter spaces of
  $\widehat{\msc}$ and of $\widehat{\msd}$ (which for simplicity we
  have denoted by $\widehat{\msf}$ too). 
\item A family of PC functors
  $\bigl\{ \msf_p: \msc_p \longrightarrow
  \msd_{\widehat{\msf}(p)}\bigr\}_{p \in
    \mathcal{P}_{\widehat{\msc}}}$, parametrized by
  $\mathcal{P}_{\widehat{\msc}}$.
\end{enumerate}
These two structures are required to have the following additional
properties. The map $\widehat{\msf}$ from~(1) above should satisfy
that for every $p,q \in \mathcal{P}_{\widehat{\msc}}$ with
$p \preceq q$ we have $\widehat{\msf}(p) \preceq
\widehat{\msf}(q)$. Moreover, we require that
$\nu(\widehat{\msf}(p)) \leq \nu(p)$ for every
$p \in \mathcal{P}_{\widehat{\msc}}$. The functors $\msf_p$ are
required to satisfy the following two conditions. First, for every
$p \preceq q$ in $\mathcal{P}_{\widehat{\msc}}$ and any choices of
comparison functors $\msh^{\widehat{\msc}}_{p,q}$,
$\msh^{\widehat{\msd}}_{\widehat{\msf}(p), \widehat{\msf}(q)}$ the
following diagram:
\begin{equation} \label{sys-fun-diag-1} \xymatrixcolsep{3pc}
  \xymatrixrowsep{3pc} \xymatrix{ \msc_p \ar[r]^-{\msf_p}
    \ar[d]_-{\msh^{\widehat{\msc}}_{p,q}} & \msd_{\widehat{\msf}(p)}
    \ar[d]^-{\msh^{\widehat{\msd}}_{\widehat{\msf}(p), \widehat{\msf}(q)}}\\
    \msc_q \ar[r]^-{\msf_q} & \msd_{\widehat{\msf}(q)}}
\end{equation}
commutes up to natural isomorphisms in the $0$-level category
$\text{pc-fun}^0(\msc_p, \msd_{\widehat{\msf}(q)})$ of the PC of
PC-functors $\text{pc-fun}(\msc_p, \msd_{\widehat{\msf}(q)})$.
Similarly, we also require that for every $p \preceq q$ in
$\mathcal{P}_{\widehat{\msc}}$ and every choice of comparison functors
$\msh^{\widehat{\msc}}_{q,p}$,
$\msh^{\widehat{\msd}}_{\widehat{\msf}(q), \widehat{\msf}(p)}$ the
diagram:
\begin{equation} \label{sys-fun-diag-2} \xymatrixcolsep{3pc}
  \xymatrixrowsep{3pc} \xymatrix{ \msc_p \ar[r]^-{\msf_p} &
    \msd_{\widehat{\msf}(p)}
    \\
    \msc_q \ar[u]^-{\msh^{\widehat{\msc}}_{p,q}} \ar[r]^-{\msf_q} &
    \msd_{\widehat{\msf}(q)}
    \ar[u]_-{\msh^{\widehat{\msd}}_{\widehat{\msf}(p),
        \widehat{\msf}(q)}} }
\end{equation}
commutes up to natural isomorphisms in the $0$-level category
$\text{pc-fun}^0(\msc_q, \msd_{\widehat{\msf}(p)})$ of the PC of
PC-functors $\text{pc-fun}(\msc_q, \msd_{\widehat{\msf}(p)})$.

It remains to define natural transformations between functors of
systems of PC's. Let $\widehat{\msc}, \widehat{\msd}$ be two systems of PC's and
$\widehat{\msf}, \widehat{\msg}: \widehat{\msc} \longrightarrow
\widehat{\msd}$ two functors having the same action on the parameter
spaces (i.e.~the two maps
$\widehat{\msf}, \widehat{\msg}: \mathcal{P}_{\widehat{\msc}}
\longrightarrow \mathcal{P}_{\widehat{\msd}}$ coincide). A natural
transformation
$\widehat{\theta}: \widehat{\msf} \longrightarrow \widehat{\msg}$ of
shift $s \in \mathbb{R}$, consists of a family of persistence natural
transformations $\theta_p : \msf_p \longrightarrow \msg_p$ for every
$p \in \mathcal{P}$ each of which of shift $s$. In addition the
$\theta_p$ are required to have the following weak compatibility with
the comparison functors. For every $p \preceq q$ and every
$X \in \Ob(\msc_p)$ there exist $0$-isomorphisms
$$\tau_{p,q}: \msh_{\widehat{\msf}(p), \widehat{\msf}(q)}(\msf_p (X))
\longrightarrow \msf_q (\msh_{p,q}(X)), \quad \sigma_{p,q}:
\msh_{\widehat{\msf}(p), \widehat{\msf}(q)}(\msg_p(X)) \longrightarrow
\msg_q(\msh_{p,q}(X)),$$ such that
$$\sigma_{p,q} \circ \msh_{\widehat{\msf}(p), \widehat{\msf}(q)}(\theta_p(X)) = 
\theta_q(\msh_{p,q}(X)) \circ \tau_{p,q}$$ in
$\hom_{\msd_{\widehat{\msf}(q)}} \bigl( \msh_{\widehat{\msf}(p),
  \widehat{\msf}(q)}(\msf_p (X)), \msg_q(\msh_{p,q}(X))
\bigr)$. Finally, we require a similar compatibility also with the
comparison functors $\msh_{q,p}$.

\subsubsection{The case of TPC's} \label{sb:sys-tpc} A system of TPC's
with increasing accuracy is simply a system $\widehat{\msc}$ of PC's
as in~\S\ref{sb:sys-pc} with the following additional
assumptions. Each of the categories $\msc_p$, $p \in \mathcal{P}$, in
the system is assumed to be a TPC. Moreover, all the functors in the
collections $\msj_{p,q}$ and $\msj_{q,p}$ are assumed to be TPC
functors. The notion of functors of PC's (and natural transformations
between them) from~\S\ref{sbsb:func-sys-pc} extends to the setting of
systems of TPC's in a straightforward way, by just requiring the
functors ${\msf}_p$ to be TPC functors.

Note also that point~\eqref{i:rem-PC-1} of Remark~\ref{rem:pc-limits}
carries over to systems
$\widehat{\msc} = \{\msc_p\}_{p \in \mathcal{P}}$ of TPC's, namely the
$\infty$-levels $\msc^{\infty}_p$, $p \in \mathcal{P}$, are mutually
equivalent in the triangulated sense.



\subsubsection{Systems of filtered
  $A_{\infty}$-categories} \label{sb:sys-ainfty} There is also a
notion of system of filtered $A_{\infty}$-categories with increasing
accuracy. The definition is analogous to that given for a system of
PC's in~\S\ref{sb:sys-pc} above, but with one significant change
regarding the comparison functors between the $q$th-category and the
$p$th one when $p \preceq q$.

In the $A_{\infty}$-case a system of categories with increasing
accuracy consists of the following. A family
$\widehat{\mathcal{A}} = \{\mathcal{A}_p\}_{p \in \mathcal{P}}$ of
filtered $A_{\infty}$-categories $\mathcal{A}_p$, parametrized by
$p \in \mathcal{P}$. Note that we do not assume these
$A_{\infty}$-categories to be pretriangulated. The parameter space
$\mathcal{P}$ is endowed with the same structures $\nu$ and $\preceq$
as in~\S\ref{sb:sys-tpc}. For $p, q \in \mathcal{P}$ with
$p \preceq q$ we have two collections $\mathcal{J}_{p,q}$ and
$\mathcal{J}_{q,p}$ of comparison $A_{\infty}$-functors with the
following properties. The functors
$\mathcal{H}_{p,q} \in \mathcal{J}_{p,q}$ from the first collection
are filtered $A_{\infty}$-functors. The functors
$\mathcal{H}_{q,p} \in \mathcal{J}_{q,p}$ from the second collection
are assumed to be $A_{\infty}$-functors
$\mathcal{H}_{q,p}: \mathcal{A}_q \longrightarrow \mathcal{A}_p$ with
linear deviation rate $\leq C | \nu(p) - \nu(q)|$, for some constant
$C$ that depends only on $\widehat{\mathcal{A}}$ (and not on $p,
q$). Moreover, the functors from $\mathcal{J}_{p,q}$ and
$\mathcal{J}_{q,p}$ are assumed to satisfy the $A_{\infty}$-analogs of
the properties listed in~\S\ref{sb:sys-tpc} on
page~\pageref{pp:sys-functors}, with the following obvious
modifications. Instead of $0$-isomorphisms $\cong_0$ between two
functors mentioned at points~\eqref{i:0-iso-sys-1}
and~\eqref{i:0-iso-sys-2} on page~\pageref{pp:sys-functors} we will
now require that we have $0$-isomorphisms between the respective functors in the persistence homology category of the
filtered $A_{\infty}$-functors
$PH (F\text{fun}_{A_{\infty}}(\mathcal{A}_p, \mathcal{A}_q))$, and
similarly for the comparison functors in
$PH(\text{fun}^{\text{LD}}(\mathcal{A}_q, \mathcal{A}_p))$. The
assumptions in point~\eqref{i:d-int-func} on
page~\pageref{pp:sys-functors} will be now replaced by requiring 
$\dint(\mathcal{F} \circ \mathcal{H}, \id) \leq C | \nu(p) - \nu(q)|$
for the interleaving distance on
$PH (\text{fun}_{A_{\infty}}(\mathcal{A}_p, \mathcal{A}_p))$, and
similarly for $\dint(\mathcal{H} \circ \mathcal{F}, \id)$.

The notion of functors between systems of PC's (as well as natural
transformations between them) has an $A_{\infty}$-analog, following
the preceding principle.

\subsubsection{Modules and Yoneda embeddings of
  systems} \label{sbsb:yoneda-sys}

Let $\widehat{\mathcal{A}} = \{\mathcal{A}_p\}_{p \in \mathcal{P}}$ be
a system of filtered $A_{\infty}$-categories. Denote by
$F\md_{\mathcal{A}_p}$ the $A_{\infty}$-category of filtered
$\mathcal{A}_p$-modules. Recall that this is a filtered
$A_{\infty}$-category (in fact a filtered dg-category). These
categories fit into a system
$$F\md(\widehat{\mathcal{A}}) := \{ F\md_{\mathcal{A}_p} \}_{p \in
  \mathcal{P}}$$ of filtered $A_{\infty}$-categories, where for
$p \preceq q$ the comparison functors
$$(\mathcal{H}_{p,q})_* : F\md_{\mathcal{A}_p} \longrightarrow
F\md_{\mathcal{A}_q}$$ are just the push-forward functors induced by
the comparison functors $\mathcal{H}_{p,q}$ of the system
$\widehat{\mathcal{A}}$, and similarly for $(\mathcal{H}_{q,p})_*$.
Since the comparison functors $\mathcal{H}_{p,q}$, for $p \preceq q$,
are filtered the same holds for $(\mathcal{H}_{p,q})_*$. As for the
functors in the opposite direction, interestingly the pushforward
functors $(\mathcal{H}_{q,p})_*$ behave somewhat better than the
$\mathcal{H}_{q,p}$ in the following sense. While $\mathcal{H}_{q,p}$
are not really filtered functors but only LD functors, the induced
comparison functors $(\mathcal{H}_{q,p})_*$ are genuinely
filtered. This is a very nice feature of the filtered push-forward
construction (see~\S\ref{ap:pshf}). However, note that it is still the
case that $(\mathcal{H}_{p,q})_* \circ (\mathcal{H}_{q,p})_*$ is only
$r$-quasi-isomorphic to $\id$, for $r = C | \nu(p) - \nu(q)|$, and the
same for $(\mathcal{H}_{q,p})_* \circ (\mathcal{H}_{p,q})_*$.


We now turn to Yoneda embeddings of systems. Let
$\widehat{\mathcal{A}}$ be a system of filtered
$A_{\infty}$-categories as above. Denote by 
$$\mathcal{Y}_p: \mathcal{A}_p \longrightarrow F\md_{\mathcal{A}_p}$$ 
the filtered Yoneda embedding (see~\S\ref{gio:fyon-def}), and let
$(\mathcal{Y}_p(\mathcal{A}_p))^{(\text{q-iso},0)}$ be the minimal
full subcategory of $F\md_{\mathcal{A}_p}$ which contains the objects
$\mathcal{Y}_p(\Ob(\mathcal{A}_p))$ and is also closed with respect to
$0$-quasi-isomorphisms. Note that the comparison functors
$(\mathcal{H}_{p,q})_*$ and $(\mathcal{H}_{q,p})_*$ preserve these
subcategories. The main point here is that for every
$A \in \Ob(\mathcal{A}_p)$ and $p \preceq q$ the modules
$(\mathcal{H}_{p,q})_* \mathcal{Y}_p(A)$ and
$\mathcal{Y}_q(\mathcal{H}_{p,q}(A))$ are
$0$-quasi-isomorphic. Similarly, for every $B \in \Ob(\mathcal{A}_q)$
the modules $(\mathcal{H}_{q,p})_*(\mathcal{Y}_q(B))$ and
$\Sigma^r \mathcal{Y}_p(\mathcal{H}_{q,p}(B))$ are
$0$-quasi-isomorphic, where $r \leq C |\nu(p) - \nu(q)|$ is the
deviation rate of the LD functor $\mathcal{H}_{q,p}$. Recall that our
categories $\mathcal{A}_p$ are endowed with a shift functor $\Sigma$,
hence the categories $\mathcal{Y}(A_p)$ and
$(\mathcal{Y}_p(\mathcal{A}_p))^{(\text{q-iso},0)}$ are both closed
under shifts.

It follows that the categories
$(\mathcal{Y}_p(\mathcal{A}_p))^{(\text{q-iso},0)}$,
$p \in \mathcal{P}$, together with the restrictions of the
push-forward comparison functors $(\mathcal{H}_{p,q})_*$,
$(\mathcal{H}_{q,p})_*$ to them, form a system of filtered
$A_{\infty}$-categories (which can be viewed as a subsystem of
$F\md(\widehat{\mathcal{A}})$). We denote this system by
$\mathcal{Y}(\widehat{\mathcal{A}})$ and call it the Yoneda system of
$\widehat{\mathcal{A}}$.

\subsubsection{Systems of PC's associated with systems of
  $A_{\infty}$-categories} \label{sbsb:sys-ai-pc} Given a system of
filtered $A_{\infty}$-categories
$\widehat{\mathcal{A}} = \{ \mathcal{A}_p\}_{p \in \mathcal{P}}$ one
can obtain a system of PC's by passing to persistence
homology. However, there is a small subtlety in this procedure. The
corresponding system $\{PH(\mathcal{A}_p)\}_{p \in \mathcal{P}}$ of
persistence homological categories is not really a system of PC's,
because for $p \preceq q$ the $A_{\infty}$-functors from
$\mathcal{J}_{q,p}$ are not filtered but only LD functors, hence do not
induce PC functors
$PH(\mathcal{A}_q) \longrightarrow PH(\mathcal{A}_p)$. However this
problem can be rectified by passing to the Yoneda system introduced
earlier in~\S\ref{sbsb:yoneda-sys}. More precisely, consider the
Yoneda system $\mathcal{Y}(\widehat{\mathcal{A}})$ of
$\widehat{\mathcal{A}}$. Recall from~\S\ref{sbsb:yoneda-sys} that this
is (genuine) system of filtered $A_{\infty}$-categories. Define
$$PH(\widehat{\mathcal{A}}) :=
\{PH((\mathcal{Y}_p(\mathcal{A}_p))^{(\text{q-iso},0)})\}_{p \in
  \mathcal{P}}$$ to be system consisting of the persistence homology
categories of $(\mathcal{Y}_p(\mathcal{A}_p))^{(\text{q-iso},0)}$
forming the Yoneda system $\mathcal{Y}(\widehat{\mathcal{A}})$. The
comparison functors $\msh_{p,q}$ are defined to be the persistence
homology functors induced by the $(\mathcal{H}_{p,q})_*$,
i.e.~$\msh_{p,q} := PH((\mathcal{H}_{p,q})_*)$ and similarly for
$\msh_{q,p}$.
  


% More precisely, denote by $\mathcal{Y}$ the Yoneda embedding functor
% and consider the $A_{\infty}$-categories
% $\mathcal{Y}(\mathcal{A}_p)^{\ziso}$, $p \in \mathcal{P}$, of filtered
% Yoneda modules (namely, the images of $\mathcal{A}_p$ under the Yoneda
% embedding followed by a completion with respect to to
% $0$-quasi-isomorphisms; see~\S\ref{sbsb:comp}).  For $p \preceq q$,
% the comparision functors $\mathcal{H}_{p,q}$ and $\mathcal{H}_{q,p}$
% induce now push forward functors
% $(\mathcal{H}_{p,q})_*: \mathcal{Y}(\mathcal{A}_p)^{\ziso}
% \longrightarrow \mathcal{Y}(\mathcal{A}_q)^{\ziso}$ and
% $(\mathcal{H}_{q,p})_*: \mathcal{Y}(\mathcal{A}_q)^{\ziso}
% \longrightarrow \mathcal{Y}(\mathcal{A}_p)^{\ziso}$. The important
% point now is that both of these functors are filtered. This is clear
% for $(\mathcal{H}_{p,q})_*$ (since $\mathcal{H}_{p,q}$ is filtered),
% but it also holds for $(\mathcal{H}_{q,p})_*$ since
% $\mathcal{H}_{q,p}$ has linear deviation.

% We now define
% $PH(\widehat{A}) := \{PH(\mathcal{Y}(\mathcal{A}_p)^{\ziso})\}_{p \in
%   \mathcal{P}}$ and obtain a system of PC's. It has the same parameter
% space $\mathcal{P}$, the same preorder $\preceq$, and size function
% $\nu$. The comparison functors
% $\msh_{p,q} = PH((\mathcal{H}_{p,q})_*)$ are the persistence homology
% functors induced from $(\mathcal{H}_{p,q})_*)$ and the same for
% $\msh_{q,p}$.

The assignment
$\widehat{\mathcal{A}} \longmapsto PH(\widehat{\mathcal{A}})$ extends
to a functor
$\text{SYS}_{F A_{\infty}} \longrightarrow \text{SYS}_{\text{PC}}$
between the category of systems of filtered $A_{\infty}$-categories
and the category of systems of PC's. This functor assigns to a functor
$\widehat{\mathcal{F}} : \widehat{\mathcal{A}} \longrightarrow \widehat{\mathcal{B}}$
between systems of filtered $A_{\infty}$-categories a functor
$PH(\widehat{\mathcal{F}}): PH(\widehat{\mathcal{A}}) \longrightarrow
PH(\widehat{\mathcal{B}})$ between the corresponding systems of PC's and this
correspondence behaves as expected on natural transformations between
such functors.

Finally, we remark that if
$\widehat{\mathcal{B}} = \{\mathcal{B}_p\}_{p \in \mathcal{P}}$ is a systems of
pre-triangulated filtered $A_{\infty}$-categories then the
corresponding system $PH(\widehat{\mathcal{B}})$ of PC's is in fact a
system of TPC's. For the rest of this paper, the most important
example in this context will be the following. Let
$\widehat{\mathcal{A}} = \{\mathcal{A}_p\}_{p \in \mathcal{P}}$ be a
system of filtered $A_{\infty}$-categories. Consider the Yoneda system
$\mathcal{Y}(\widehat{\mathcal{A}})$ as defined
in~\S\ref{sbsb:yoneda-sys}. By applying the construction
from~\S\ref{sbsb:fmod} to each category that participates in the
latter system we obtain a new system
$\mathcal{Y}(\mathcal{\widehat{A}})^{\Delta} =
\{\mathcal{Y}(\mathcal{A}_p)^{\Delta} \}_{p \in \mathcal{P}}$ of
filtered pre-triangulated categories. The comparison functors for this
system are induced from those of the Yoneda system
(see~\S\ref{sbsb:yoneda-sys}). This works since filtered and LD
$A_{\infty}$-functors preserve triangulated completions as defined
in~\S\ref{sbsb:fmod}. We now pass to persistence homology and define
the following system of TPC's:
\begin{equation} \label{eq:PDA-sys} PD(\widehat{\mathcal{A}}) :=
  \{PH(\mathcal{Y}(\mathcal{A}_p)^{\Delta})\}_{p \in \mathcal{P}}.
\end{equation}
We call this system the {\em persistence derived system of}
$\widehat{\mathcal{A}}$.


\subsubsection{Systems with a coherent base of
  objects} \label{sbsb:sys-base} Let
$\widehat{\msc} = \{\msc_p\}_{p \in \mathcal{P}}$ be a system of PC's
with increasing accuracy. A coherent base of objects $\{mathscr\}{B}$ for
$\widehat{\msc}$ is a family of subsets
  $\mathscr{B} = \{\mathscr{B}_p\}_{p \in \mathcal{P}}$ of objects
  $\mathscr{B}_p \subset \Ob(\msc_p)$ for every $p \in \mathcal{P}$,
satisfying the following properties:
\begin{enumerate}
\item For every $p \preceq q$, every two comparison functors
  $\msh'_{p,q}, \msh''_{p,q} \in \mathscr{J}_{p,q}$ and every
  $A \in \mathscr{B}_p$ we have $\msh'_{p,q}(A) = \msh''_{p,q}(A)$ and
  $\msh'_{p,q}(A) \in \mathscr{B}_q$. We require the analogous
  properties to hold also for every two comparison functors
  $\msh'_{q,p}, \msh''_{q,p} \in \mathscr{J}_{q,p}$ and every
  $B \in \mathscr{B}_q$.
\item For every $p \preceq q \preceq r$, every two comparison functors
  $\msh_{p,q} \in \mathscr{J}_{p,q}$,
  $\msh_{q,r} \in \mathscr{J}_{q,r}$ and every $A \in \mathscr{B}_p$
  we have $\msh_{q,r}(\msh_{p,q}(A)) = \msh_{p,r}(A)$. We require the
  analogous properties to hold also for all comparison functors going
  in the other direction, namely for $\msh_{r,q}$, $\msh_{q,p}$,
  $\msh_{r,p}$ and
  every object $B \in \mathscr{B}_r$.
\item For every $p \preceq q$, every choice of comparison functors
  $\msh_{p,q}$, $\msh_{q,p}$ and every $A \in \mathscr{B}_p$,
  $B \in \mathscr{B}_q$ we have $\msh_{q,p}(\msh_{p,q}(A)) = A$,
  $\msh_{p,q}(\msh_{q,p}(B)) = B$.
\end{enumerate}
It follows from the definition above that for every
$p', p'' \in \mathcal{P}$ we have a canonical bijection
$\mathscr{B}_{p'} \longrightarrow \mathcal{B}_{p''}$, namely
$\msh_{q, p''} \circ \msh_{p', q}$ where $q \in \mathcal{P}$ is any
element with $q \succeq p', p''$.

One can define functors between systems of PC's with a coherent base
of objects in a straightforward way.

The concept of systems with a coherent base of objects (as well
functors between such) carries over without any modification also to
the case of systems of filtered $A_{\infty}$-categories.

A special case of a system of categories (filtered $A_{\infty}$ or
PC's) with a coherent base of objects that will appear later is when
$\mathscr{B}_p = \Ob(\msc_p)$ for every $p \in \mathcal{P}$. We will
refer to such a situation as a {\em coherent full base of objects}.
In the following we will encounter an even simpler situation, where
all the categories $\msc_p$, $p \in \mathcal{P}$, have the same set of
objects (i.e.~$\Ob(\mathcal{C}_{p'}) = \Ob(\mathcal{C}_{p''})$ for
every $p', p'' \in \mathcal{P}$) and all the comparison functors act
like the identity on objects, namely $\msh_{p,q}(X) = X$ for every
$p \preceq q$, every comparison functor $\msh_{p,q}$, and
$X \in \Ob(\mathcal{C}_{p})$. In this case we will view the coherent
base $\mathscr{B}$ as a single set rather than a family, and set
$\mathscr{B} = \Ob(\mathcal{C}_p)$ for any $p \in \mathcal{P}$.  We
will call this type of systems by the name {\em system of categories
  (PC's or filtered $A_{\infty}$-categories) with a fixed full base of
  objects}.

\begin{rem} \label{r:coh-base} Consider a system of filtered
  $A_{\infty}$-categories
  $\widehat{\mathcal{A}} = \{\mathcal{A}_p\}_{p \in \mathcal{P}}$ with
  a fixed full base of objects $\mathcal{B}$. Let
  $F\md(\widehat{\mathcal{A}})$ be the corresponding system of
  filtered modules over the $\mathcal{A}_p$'s and
  $\mathcal{Y}(\widehat{\mathcal{A}})$ be the Yoneda system of
  $\widehat{\mathcal{A}}$ (see~\S\ref{sbsb:yoneda-sys}). Unlike
  $\widehat{\mathcal{A}}$ these other two systems do not have a fixed
  (or even coherent) full base of objects. The reason is that the
  comparison functors for $F\md(\widehat{\mathcal{A}})$ and
  $\mathcal{Y}(\widehat{\mathcal{A}})$ are the ``module push-forward''
  functors $\mathcal{H}_*$ induced from the comparison functors
  $\mathcal{H}$ of $\widehat{\mathcal{A}}$ and these do not satisfy
  the conditions required for a coherent base of objects. More
  specifically, for $p \preceq q \preceq r$, the functors
  $(\mathcal{H}_{q,r})_* \circ (\mathcal{H}_{p,q})_*$ and
  $(\mathcal{H}_{q,r})_*$ do not act in the same way on filtered
  modules (though the images of a module under these two functors are
  $0$-quasi-isomorphic). Similarly, the push forward
  $(\mathcal{H}_{p,q})_* \mathcal{Y}(X)$ of a Yoneda module is not a
  Yoneda module but only $0$-quasi-isomorphic to a Yoneda module.
\end{rem}


\subsubsection{Homotopy systems} \label{sbsb:homotopy-sys}

Let $\widehat{\mathcal{A}} = \{\mathcal{A}_p\}_{p \in \mathcal{P}}$ be
a system of filtered $A_{\infty}$-categories with a coherent full base
of objects. We say that the system $\widehat{\mathcal{A}}$ is a {\em
  homotopy system} if the following conditions hold:
\begin{enumerate}
\item For every $p \preceq q$ in $\mathcal{P}$, and every two
  comparison functors
  $\mathcal{H}'_{p,q}, \mathcal{H}''_{p,q} \in \mathcal{J}_{p,q}$
  there is a homotopy
  $$T \in \hom_{F \text{fun}(\mathcal{A}_p,
    \mathcal{A}_q)}^0(\mathcal{H}'_{p,q}, \mathcal{H}''_{p,q})$$
  between $\mathcal{H}'_{p,q}$ and $\mathcal{H}''_{p,q}$ that does not
  shift filtrations. From now on we will call such $T$'s
  $0$-homotopies and say that $\mathcal{H}'_{p,q}$ and
  $\mathcal{H}''_{p,q}$ are $0$-homotopic, and the superscript in
  $\hom^0$ stands for pre-natural transformations of filtration level
  $0$ (i.e.~they preserve filtrations). We refer
  to~\cite[Section~(1h)]{Se:book-fukaya-categ} for the definition of
  homotopy between $A_{\infty}$-functors in the unfiltered case.
\item For every $p \preceq q$ in $\mathcal{P}$,  every two
  comparison functors
  $\mathcal{H}'_{q,p}, \mathcal{H}''_{q,p} \in \mathcal{J}_{q,p}$ are
  $0$-homotopic as LD-functors with deviation rate
  $C|\nu(p) - \nu(q)|$, where $C$ is the constant associated with the
  family $\widehat{\mathcal{A}}$ that appears in the definition of
  systems of filtered $A_{\infty}$-categories
  in~\S\ref{sb:sys-ainfty}. Specifically, this means that there exists
  a homotopy $T$ between the LD-functors
  $(\mathcal{H}'_{q,p}, C|\nu(p) - \nu(q)|)$ and
  $(\mathcal{H}''_{q,p}, C|\nu(p) - \nu(q)|)$ of shift $0$, namely its
  $d$th order $T_d$ shifts filtrations by $\leq d C|\nu(p) -
  \nu(q)|$. See~\S\ref{a:func-LD}.
\item For every $p \preceq q \preceq r$ in $\mathcal{P}$ and every
  choices of comparison functors $\mathcal{H}_{p,q}$,
  $\mathcal{H}_{q,r}$, $\mathcal{H}_{p,r}$ we have that
  $\mathcal{H}_{q,r} \circ \mathcal{H}_{p,q}$ and $\mathcal{H}_{p,r}$
  are $0$-homotopic.
\end{enumerate}

\subsubsection{Invariants of systems} \label{sb:sys-inv} Systems of
categories $\widehat{\msc} = \{\msc_p\}_{p \in \mathcal{P}}$ with
increasing accuracy call for considering their limits (or rather
colimits) when the parameter $p \in \mathcal{P}$ goes to ``infinity''
(or equivalently when the accuracy-levels $\nu(p)$ converges to
$0$). In other words, it would be desirable to be able to define a
colimit category $\colim_{p \in \mathcal{P}} \msc_p$ which will have
the same structures (e.g.~PC or TPC structures, filtered
$A_{\infty}$-structure etc.) as the categories $\msc_p$ forming the
family $\widehat{\msc}$. This seems a non-trivial foundational problem
and we will not attempt to solve it in this paper. Instead, we will
show that some algebraic invariants and measurements associated with
the categories $\msc_p$ do admit colimits, even without having a
colimit category. We will concentrate on three such invariants (each
of which is defined in a different settings): the interleaving
(pseudo)-distance on the set of objects of a PC, the Grothendieck
groups $K_0$ of triangulated categories and Hochschild homologies of
$A_{\infty}$-categories.

We begin with a useful equivalence relation on the totality of
  objects of a system of categories.  \label{p:equiv-ob} Let
  $\widehat{\msc} = \{\msc_p\}_{p \in \mathcal{P}}$ be a system of
  PC's with increasing accuracy. Consider
  $$\Ob^{\text{tot}}(\widehat{\msc}) := \coprod_{p\in \mathcal{P}}
  \Ob(\msc_p)$$ and denote elements of this set by $(A,p)$ with
  $p \in \mathcal{P}$ and $A \in \msc_p$. Define an equivalence
  relation $\sim_{\widehat{\msc}}$ on
  $\Ob^{\text{tot}}(\widehat{\msc})$ as follows:
  $(A', p') \sim_{\widehat{\msc}} (A'', p'')$ if there exists
  $q \in \mathcal{P}$ with $p',p'' \preceq q$ such that
  $\msh_{p',q}(A') \cong_0 \msh_{p'',q}(A'')$ for some comparison
  functors $\msh_{p',q} \in \msj_{p',q}$ and
  $\msh_{p'',q} \in \msj_{p'',q}$. Here $\cong_0$ means an isomorphism
  in $\msc_q^0$. (Since all the functors in each of $\msj_{p',q}$,
  $\msj_{p'',q}$, are mutually $0$-isomorphic, requiring that
  $\msh_{p',q}(A') \cong_0 \msh_{p'',q}(A'')$ for some
  $\msh_{p',q} \in \msj_{p',q}$, $\msh_{p'',q} \in \msj_{p'',q}$, is
  the same as to ask that the same property holds for {\em all}
  $\msh_{p',q} \in \msj_{p',q}$, $\msh_{p'',q} \in \msj_{p'',q}$.)
  Clearly if $A', A'' \in \Ob(\msc_p)$ are $0$-isomorphic then
  $(A', p) \sim_{\widehat{\msc}} (A'',p)$, but in general the
  equivalence relation $\sim_{\widehat{\msc}}$ may relate more objects
  belonging to categories parametrized by different $p$'s.
Denote by $\widetilde{\Ob}(\widehat{\msc})$ the equivalence
  classes of $\sim_{\widehat{\msc}}$.  Given
  $\mathcal{S} \subset \widetilde{\Ob}(\widehat{\msc})$ and
  $p \in \mathcal{P}$ we write
$$\mathcal{S}_p := \{ A \in \Ob(\msc_p) \mid [A]_{\widehat{\msc}}
\in \mathcal{S}\},$$ where $[A]_{\widehat{\msc}}$ stands for the
equivalence class of $A$ in $\widetilde{\Ob}(\widehat{\msc})$.  We
will use the same notation $\mathcal{S}_p$ also when we deal with just
one element $\mathcal{S} \in \widetilde{\Ob}(\widehat{\msc})$.

We will define now a pseudo distance on
  $\widetilde{\Ob}(\widehat{\msc})$ induced by the interleaving
  distances on the $\msc_p$'s.
  Denote by $d^p_{\text{int}}$
  the interleaving (pseudo) distance on $\Ob(\msc_p)$. We define the
  following measurement on pairs of elements
  $\widetilde{X}, \widetilde{Y} \in \widetilde{\Ob}(\widehat{\msc})$:
\begin{equation} \label{eq:dist-lim}
  \widetilde{d}_{\text{int}}(\widetilde{X}, \widetilde{Y}) = \inf
  \bigl\{ d^p_{\text{int}}(X',Y') \mid p \in \mathcal{P}, X' \in
  \widetilde{X}_p, Y' \in \widetilde{Y}_p \bigr\},
\end{equation}
and call it the {\em limit interleaving distance} (the wording "limit"
will be justified shortly, in Lemma~\ref{l:lim-dint} below). Note that
the term $d^p_{\text{int}}(X',Y')$ that appears in~\eqref{eq:dist-lim}
depends only on the $0$-isomorphism classes of
$X', Y' \in \Ob(\msc_p)$. It is straightforward to see that
$\widetilde{d}_{\text{int}}$ is a pseudo-distance. Note that due to
the $\inf$ in~\eqref{eq:dist-lim},
$\widetilde{d}_{\text{int}}(\widetilde{X}, \widetilde{Y})$ may vanish
even if $\widetilde{X} \neq \widetilde{Y}$. It may also assume the
value $\infty$ in case $\widetilde{X}$ and $\widetilde{Y}$ do not
represent isomorphic objects in the persistence $\infty$-level
categories $\msc^{\infty}_p$ for ``large'' $p$'s. The pseudo-metric
$\widetilde{d}_{\text{int}}$ can be viewed as a more robust version of
the pseudo-distance $\hat{d}$ from Remark
\ref{rem:def-TPC-approx}. Indeed, as its name suggests, the limit
distance can also be defined in terms of the limit of
$d^p_{\text{int}}$ as $\nu(p) \longrightarrow 0$.
More
specifically:
\begin{lem} \label{l:lim-dint} Let $\{p_k\}_{k \in \mathbb{N}}$ be a
  sequence of elements from $\mathcal{P}$ with
  $\nu(p_k) \longrightarrow 0$ as $k \longrightarrow \infty$. Then for
  every
  $\widetilde{X}, \widetilde{Y} \in \widetilde{\Ob}(\widehat{\msc})$
  we have
  $$\lim_{k \to \infty} d^{p_k}_{\text{int}}(X'_k,Y'_k) =
  \widetilde{d}_{\text{int}}(\widetilde{X},
  \widetilde{Y}),$$ where $X'_k, Y'_k$ are any sequences with $X'_k
  \in \widetilde{X}_{p_k}$, $Y'_k \in \widetilde{Y}_{p_k}$.
\end{lem}

\begin{proof}
  This follows from elementary arguments, using the properties of the
  comparison functors listed on page~\pageref{pp:sys-functors}
  together with Lemma~\ref{l:d-func}.
\end{proof}


  \begin{rem} \label{r:d-limit} Let
    $\widehat{\msf}: \widehat{\msc} \longrightarrow \widehat{\msd}$ be
    a functor between two systems of PC's
    (see~\S\ref{sbsb:func-sys-pc}).  It follows directly from the
    definitions that $\widehat{\msf}$ induces a well defined map
    $\widetilde{\msf}: \widetilde{\Ob}(\widehat{\msc})
    \longrightarrow\widetilde{\Ob}(\widehat{\msd})$ that satisfies
    $\widetilde{\msf}([A]_{\widehat{\msc}}) =
    [\msf_p(A)]_{\widehat{\msd}}$ for every $A \in \Ob(\msc_p)$,
    $p \in \mathcal{P}$. Moreover, $\widetilde{\msf}$ is non-expanding
    with respect to the limiting interleaving distances on
    $\widetilde{\Ob}(\widehat{\msc})$ and
    $\widetilde{\Ob}(\widehat{\msd})$ in the sense that
    $\widetilde{d}^{\widehat{\msd}}_{\text{int}}
    (\widetilde{\msf}(\widetilde{X}), \widetilde{\msf}(\widetilde{Y}))
    \leq \widetilde{d}^{\widehat{\msc}}_{\text{int}}(\widetilde{X},
    \widetilde{Y})$ for every
    $\widetilde{X}, \widetilde{Y} \in
    \widetilde{\Ob}(\widehat{\msc})$.
\end{rem}


We now turn to $K$-groups. Let
$\widehat{\msc} = \{\msc_p\}_{p \in \mathcal{P}}$ be a system of TPC's
(see~\S\ref{sb:sys-tpc}). Recall that the $0$-persistence level
$\msc_p^0$ of each $\msc_p$ is a triangulated category. Denote by
$K_0(\msc_p^0)$ the Grothendieck (or $K$-) group of $\msc_p^0$. Let
$p \preceq q$. Recall that all the comparison functors
$\msh_{p,q}: \msc_p \longrightarrow \msc_q$ are TPC-functors, hence
they restrict to triangulated functors
$\msh^0_{p,q}: \msc^0_p \longrightarrow \msc^0_q$. We thus obtain an
induced homomorphism
\begin{equation} \label{eq:HK-pq} \mathscr{H}^K_{p,q}: K_0(\msc^0_p)
  \longrightarrow K_0(\msc^0_q).
\end{equation}
This homomorphism turns out to be independent of the specific choice
of the comparison functor $\mathcal{H}_{p,q}$. To see this, recall
that for $p \preceq q$ every two comparison functors
$\msh'_{p,q}, \msh''_{p,q} \in \mathscr{J}_{p,q}$
are $0$-isomorphic, and
that if two objects $X, Y \in \Ob(\mathscr{E})$ in a triangulated
category $\mathscr{E}$ are isomorphic, then their classes in
$K_0(\mathscr{E})$ are equal.

A similar argument shows that the $K$-group homomorphisms
from~\eqref{eq:HK-pq} satisfy
$\msh^K_{q,r} \circ \msh^K_{p,q} = \msh^K_{p,r}$ for every
$p\preceq q \preceq r$, and clearly $\msh^K_{p,p} = \id$ for every
$p$. It follows that $\{K_0(\msc_p^0)\}_{p \in \mathcal{P}}$ together
with the maps $\msh^{K}_{p,q}$ form a directed system of abelian
groups.  We define the $K_0$-group of $\widehat{\msc}$ to be the
colimit of this system:
\begin{equation} \label{eq:K-colim} K_0(\widehat{\msc}^0) := \colim_{p
    \in \mathcal{P}} K_0(\msc_p^0).
\end{equation}
The assignment $\widehat{\msc} \longmapsto K_0(\widehat{\msc}^0)$ is
functorial in the sense that TPC-functors
$\widehat{\mathscr{F}}: \widehat{\msc} \longrightarrow \widehat{\msd}$
between systems of TPC's
(see~\S\ref{sbsb:func-sys-pc},~\S\ref{sb:sys-tpc}) induce
homomorphisms
$\widehat{\mathscr{F}}^K: K_0(\widehat{\msc}^0) \longrightarrow
K_0(\widehat{\msd}^0)$.

\begin{rem} \label{r:K-Nov} Denote by $\Lambda^P_{\mathbb{Z}}$ the
  ring of Novikov polynomials (i.e.~finite Novikov series) with
  coefficients in $\mathbb{Z}$. Its elements can be written as finite
  sums $\sum n_k t^{a_k}$ with $n_k \in \mathbb{Z}$,
  $a_k \in \mathbb{R}$ ($t$ stands for the formal Novikov
  variable). As explained in~\cite{BCZ:PKth}, the Grothendieck group
  $K_0(\msc^0)$ of the $0$-persistence level category $\msc^0$ of a
  TPC $\msc$ is a $\Lambda^P_{\mathbb{Z}}$-module, where the action of
  $t^a$ on the class $[X]$ of an object $X \in \Ob(\msc^0)$ is defined
  by $t^{a}[X] := [\Sigma^a X]$ (here $\Sigma$ is the shift functor of
  $\msc$). It is straightforward to check that the homomorphisms
  $\msh^K_{p,q}$ are $\Lambda^P_{\mathbb{Z}}$-linear and therefore the
  colimit $K_0(\widehat{\msc})$ inherits the structure of a
  $\Lambda^P_{\mathbb{Z}}$-module. Similarly, the homomorphisms
  $\widehat{\mathscr{F}}^K$ induced by functors $\widehat{\msf}$ of TPC-systems
  are maps of $\Lambda^P_{\mathbb{Z}}$-modules.
\end{rem}

Finally, we discuss persistence Hochschild invariants of
  systems. Let
$\widehat{\mathcal{A}} = \{\mathcal{A}_p\}_{p \in \mathcal{P}}$ be a
systems of $A_{\infty}$-categories with increasing accuracy, and with a
coherent full base of objects (see~\S\ref{sbsb:sys-base}). Assume
further that $\widehat{\mathcal{A}}$ is a homotopy system as defined
in~\S\ref{sbsb:homotopy-sys}.

For every $p \in \mathcal{P}$ we have the persistence Hochschild
homology $PHH(\mathcal{A}_p, \mathcal{A}_p)$ of $\mathcal{A}_p$ (with
coefficients in the diagonal bimodule of $\mathcal{A}_p$). Let
$p, q \in \mathcal{P}$ with $p \preceq q$. Each comparison functor
$\mathcal{H}_{p,q} : \mathcal{A}_p \longrightarrow \mathcal{A}_q$
induces a map of persistence modules
$$\mathcal{H}^{\text{PHH}}_{p,q}: PHH(\mathcal{A}_p, \mathcal{A}_p)
\longrightarrow PHH(\mathcal{A}_q, \mathcal{A}_q).$$ Since every two
comparison functors from $\mathcal{J}_{p,q}$ are homotopic it follows
from Proposition~\ref{p:homotopy-HH} that
$\mathcal{H}^{\text{PHH}}_{p,q}$ is independent of the specific choice
of $\mathcal{H}_{p,q}$. Moreover, for every $p \preceq q \preceq r$ we
have that
$\mathcal{H}^{\text{PHH}}_{q,r} \circ \mathcal{H}^{\text{PHH}}_{p,q} =
\mathcal{H}^{\text{PHH}}_{p,r}$, and
$\mathcal{H}^{\text{PHH}}_{p,p} = \id$ for every $p \in \mathcal{P}$.
In other words,
$\{PHH(\mathcal{A}_p, \mathcal{A}_p)\}_{p \in \mathcal{P}}$ together
with the maps $\mathcal{H}^{\text{PHH}}_{p,q}$, $p \preceq q$, forms a
directed system of persistence modules, parametrized by $\mathcal{P}$.
We define the persistence Hochschild homology of
$\widehat{\mathcal{A}}$ to be the colimit of this system:
\begin{equation} \label{eq:HH-colim} PHH(\widehat{\mathcal{A}}) :=
  \colim_{p \in \mathcal{P}} PHH(\mathcal{A}_p, \mathcal{A}_p).
\end{equation}
Note that since colimits of directed systems of persistence modules
are persistence modules, $PHH(\widehat{\mathcal{A}})$ is a persistence
module too.

\begin{rem} \label{r:K-HH} Let $\mathcal{A}$ be a strictly unital
  $A_{\infty}$-category and denote its derived category by
  $D(\mathcal{A})$. There is a canonical map
  \begin{equation} \label{eq:K-HH-1} \kappa: K_0(D(\mathcal{A}))
    \longrightarrow HH_*(\mathcal{A}, \mathcal{A})
  \end{equation}
  which satisfies $\kappa([X]) = [e_X]$ for every
  $X \in \Ob(\mathcal{A})$. Here $[X] \in K_0(D(\mathcal{A}))$ stands
  for the $K_0$-class of $X$, $e_X \in \hom_{\mathcal{A}}(X,X)$ is the
  unit and $[e_X] \in HH(\mathcal{A}, \mathcal{A})$ is its Hochschild
  homology class. The map $\kappa$ is the composition of two maps,
  $\kappa = j^{-1} \circ \kappa'$, where
  $\kappa' : K_0(D(\mathcal{A})) \longrightarrow
  HH(\mathcal{Y}(\mathcal{A})^{\Delta},
  \mathcal{Y}(\mathcal{A})^{\Delta})$ and
  $j: HH(\mathcal{A}, \mathcal{A}) \longrightarrow
  HH(\mathcal{Y}(\mathcal{A})^{\Delta},
  \mathcal{Y}(\mathcal{A})^{\Delta})$. The map $\kappa'$ is defined by
  $\kappa([X]) = [e_X]$ for every $X \in \Ob(D(\mathcal{A}))$ and its
  well-definedness follows from simple considerations (e.g.~by
  using~\cite[Proposition~3.8]{Se:book-fukaya-categ}). The map $j$ is
  the one induced in homology by the chain level inclusion of the
  Hochschild complex of $\mathcal{A}$ into that of
  $\mathcal{Y}(\mathcal{A})^{\Delta}$. A version of Morita invariance
  for Hochschild homology states that $j$ is an isomorphism, hence we
  can define $\kappa = j^{-1} \circ \kappa'$. (For Morita invariance
  of Hochschild homology see e.g.~\cite{To:dg-cat} in the case when
  $\mathcal{A}$ is a dg-category, and the discussion
  in~\cite[Section~5]{Se:homology-spheres}; the case of $A_{\infty}$-categories is treated in~\cite{Sher:form-hodge}. See
  also~\cite[Section~5.4]{Bi-Co:lefcob-pub} for the significance of
  the map~\eqref{eq:K-HH-1} in some geometric situations.)

  We expect the same considerations to continue to hold also in our
  persistence setting. More specifically, let
  $\widehat{\mathcal{A}} = \{\mathcal{A}_p\}_{p \in \mathcal{P}}$ be a
  system of filtered $A_{\infty}$-categories. We expect to have for
  every $p \in \mathcal{P}$ a homomorphism
  $\kappa_p : K_0(PD(\mathcal{A}_p)^0) \longrightarrow
  PHH^0(\mathcal{A}_p, \mathcal{A}_p)$, where the $0$-superscript on
  the left stands for the $0$-persistence level category of
  $PD(\mathcal{A}_p)$ and the one on the right denotes the
  $0$-persistence level of the persistence module
  $PHH(\mathcal{A}_p, \mathcal{A}_p)$. The maps $\kappa_p$ are then
  expected to be compatible with the respective directed systems
  parametrized by $p \in \mathcal{P}$ and to induce one map
  $$\widehat{\kappa}: K_0(PD(\widehat{\mathcal{A}})^0) \longrightarrow
  PHH^0(\widehat{\mathcal{A}}).$$
\end{rem}


\subsubsection{Approximability revisited} \label{s:approximability}
%\subsubsection{Approximability within a system} \label{sb:approx-sys}

Let $\widehat{\msc}$ be a system of TPC's with increasing accuracy, as
defined above. We will use here the equivalence relation
  $\sim_{\widehat{\msc}}$ on
  $\Ob^{\text{tot}}(\widehat{\msc}) = \coprod_{p\in \mathcal{P}}
  \Ob(\msc_p)$ as defined in~\S\ref{sb:sys-inv} on
  page~\pageref{p:equiv-ob} and the notation from that section.
\begin{dfn} \label{d:approx-sys} Let
  $\mathcal{L} \subset \widetilde{\Ob}(\widehat{\msc})$ and let
  $\mathcal{F}_1, \ldots, \mathcal{F}_l \in
  \widetilde{\Ob}(\widehat{\msc})$ be a finite collection (of equivalence
  classes of objects), and let $\epsilon>0$. We say that $\mathcal{L}$
  is $\epsilon$-approximable by
  $\{\mathcal{F}_1, \ldots, \mathcal{F}_l\}$ if there exists
  $\delta>0$ and $c_{\epsilon}>0$ such that for every
  $p \in \mathcal{P}$ with $\nu(p) < \delta$ we have
  $${d}_{\text{int}} \Bigl(\mathcal{L}_p,
  \Ob \langle \{F_1, \ldots, F_l\}_p \rangle^{\Delta} \Bigr)
  < \epsilon + c_{\epsilon}\nu(p).$$
  Let $\mathcal{F} \subset \widetilde{\Ob}(\widehat{\msc})$. We say that
  $\mathcal{L}$ is approximable by $\mathcal{F}$ if for every
  $\epsilon>0$ there exists a finite collection
  $\mathcal{F}_1, \ldots, \mathcal{F}_l \in \mathcal{F}$ which
  $\epsilon$-approximates $\mathcal{L}$ and the constants $c_{\epsilon}$ 
  are bounded, independently of $\epsilon$.

  \end{dfn}
  Recall that 
  %$\{F_1, \ldots, F_l\}_p^{\Sigma}$ stands
  %for completing the set of objects
  %$\{F_1, \ldots, F_l\}_p \subset \Ob(\msc_p)$ with respect to all
  %shifts, and 
  $\langle \mathcal{O} \rangle^{\Delta}$ stands for the
  smallest sub-TPC of $\msc_p$ containing the set of
  objects $\mathcal{O}$. Sometimes we will use the wording {\em
    $\{\mathcal{F}_1, \ldots, \mathcal{F}_l\}$ $\epsilon$-approximates
    $\mathcal{L}$} and {\em $\mathcal{F}$ approximates $\mathcal{L}$}.

\

Similarly we also define retract approximability as follows. 
\begin{dfn} \label{d:approx-sys-ret} Let $\widehat{\msc}$,
  $\mathcal{L}$, $\mathcal{F}_1, \ldots, \mathcal{F}_l$ be as in
  Definition~\ref{d:approx-sys} and let $\epsilon>0$. We say that
  $\mathcal{L}$ is $\epsilon$-retract-approximable by
  $\{\mathcal{F}_1, \ldots, \mathcal{F}_l\}$ if there exists
  $\delta>0$, $c_{\epsilon}>0$, such that for every $p \in \mathcal{P}$ with
  $\nu(p) < \delta$ we have
  $${d}_{\text{r-int}} \Bigl(\mathcal{L}_p,
  \Ob \langle \{F_1, \ldots, F_l\}_p \rangle^{\Delta} \Bigr)
  < \epsilon +c_{\epsilon}\nu(p).$$

  Let $\mathcal{F} \subset \widetilde{\Ob}(\widehat{\msc})$. We say
  that $\mathcal{L}$ is retract-approximable by $\mathcal{F}$ if for
  every $\epsilon>0$ there exists a finite collection
  $\mathcal{F}_1, \ldots, \mathcal{F}_l \in \mathcal{F}$ which
  $\epsilon$-retract-approximates $\mathcal{L}$ and the constants
  $c_{\epsilon}$ are uniformly bounded.
\end{dfn}

\begin{rem}\label{rem:different_approx} It is useful at this point to
  relate our two notions of approximability, in Definitions
  \ref{def:TPC-approx} and \ref{d:approx-sys}.  In our geometric
  applications, the set $\mathcal{L}$ will be a set of Lagrangian
  submanifolds $\mathcal{L}ag(M)$ and the category $\msc_{p}$ will be
  a certain TPC associated with a filtered Fukaya category defined using
  perturbation data that is identified by the parameter $p$.  The
  inclusion $\mathcal{L}ag(M)\subset \msc_{p}$ is the Yoneda
  embedding. The set $\mathcal{L}ag(M)$ is endowed with the spectral
  metric $d_{\gamma}$ (that will be recalled later) and this metric is
  closely related to the stabilized interleaving metric $\sdint$ from
  \eqref{eq:sdint} restricted to $\mathcal{L}ag(M)$.  In view of this,
  we will see that $\epsilon$ - approximability in the sense of
  Definition \ref{d:approx-sys} implies $\epsilon'$-approximability
  for all $\epsilon'>\epsilon$ in the sense of Definition
  \ref{def:TPC-approx}. The reason for this discrepancy in the
  parameters $\epsilon,\epsilon'$ is related to the presence of the
  term $c_{\epsilon}\nu(p)$ in Definition \ref{d:approx-sys}.  Given
  that these properties are supposed to be satisfied for all
  $\epsilon,\epsilon' >0$ and $\nu(p)$ can be taken as small as
  desired, this discrepancy is essentially irrelevant. However, for a
  fixed $\epsilon$, including the term $c_{\epsilon} \nu(p)$ in the
  upper bound in Definition \ref{d:approx-sys} is useful for questions
  having to do with minimizing the cardinality of the set
  $\{\mathcal{F}_{i}\}$ for that fixed $\epsilon$.  A similar remark
  applies to retract-approximability.
\end{rem}

\subsubsection{Stabilizing functors} \label{sbsb:sfunc} Let
  $\widehat{\msc} = \{ \msc_p\}_{p \in \mathcal{P}}$ be a system of
  PC's with increasing accuracy and let
  $\widehat{\msf}: \widehat{\msc} \longrightarrow \widehat{\msc}$ be
  an endo-functor of such systems (see~\S\ref{sbsb:func-sys-pc}). Let
  $\mathscr{B} \subset \Ob^{\text{tot}}(\widehat{\msc})$. We say that
  $\widehat{\msf}$ is {\em $\mathscr{B}$-stabilizing} if the following
  holds for every $p \in \mathcal{P}$:
  \begin{enumerate}
  \item \label{i:stab-1} For every $A \in \mathscr{B}_p$ we have
    $(A,p) \sim_{\widehat{\msc}} (\msf_p(A), \widehat{\msf}(p))$,
    where $\sim_{\widehat{\msc}}$ is the equivalence relation
    from~\S\ref{sb:sys-inv}, page~\pageref{p:equiv-ob}.
  \item \label{i:stab-2} For every $A', A'' \in \mathscr{B}_p$ there
    exists $q \in \mathcal{P}$ with $p, \widehat{\msf}(p) \preceq q$
    and two $0$-isomorphisms
    $\sigma_{A'}: \msh_{\widehat{\msf}(p), q} \msf_p A'
    \longrightarrow \msh_{p,q}A'$,
    $\sigma_{A''}: \msh_{\widehat{\msf}(p), q} \msf_p A''
    \longrightarrow \msh_{p,q}A''$ such that for every
    $u \in \hom_{\msc_p}(A',A'')$ we have:
    \begin{equation} \label{eq:stab-2} \sigma_{A''} \circ
      \msh_{\widehat{\msf}(p), q} \circ \msf_p(u) = \msh_{p,q}(u)
      \circ \sigma_{A'}.
    \end{equation}
  \end{enumerate}

\begin{rem} \label{r:stab-func} 
Condition~\eqref{i:stab-1}
    above is equivalent to requiring that the map
    $\widetilde{\msf}: \widetilde{\Ob}(\widehat{\msc}) \longrightarrow
    \widetilde{\Ob}(\widehat{\msc})$ (see Remark~\ref{r:d-limit})
    stabilizes $\mathscr{B}$ (i.e.~$\widetilde{\msf}$ sends
    $\mathscr{B}$ to itself and restricts to the identity map
    $\mathscr{B} \longrightarrow \mathscr{B}$).

  Condition~\eqref{i:stab-2} can be rephrased as follows. By a
    slight abuse of notation denote for every $p \in \mathcal{P}$ by
    $\mathscr{B}_p \subset \msc_p$ the full persistence subcategory
    with objects $\mathscr{B}_p$. We require that given
    $p \in \mathcal{P}$, for large enough (with respect to $\preceq$)
    $q \in \mathcal{P}$ the two restricted PC functors
    $$\msh_{\widehat{\msf}(p), q} \circ \msf_p|_{\mathscr{B}_p}, \, 
    \msh_{p,q}|_{\mathscr{B}_p}: \mathscr{B}_p \longrightarrow
    \mathscr{B}_q$$ are naturally $0$-isomorphic in the category of
    persistence functors
    $\mathscr{B}_p \longrightarrow \mathscr{B}_q$.
  \end{rem}

  In the presence of approximating families in TPC's,
    stabilizing functors have a remarkable metric property.
    \begin{prop} \label{p:stab-func} Let $\widehat{\msc}$ be a system
      of TPC's, $\mathcal{L} \subset \widetilde{\Ob}(\widehat{\msc})$
      and
      $\mathcal{F}_1, \ldots, \mathcal{F}_l \in
      \widetilde{\Ob}(\widehat{\msc})$ an $\epsilon$-approximating
      family for $\mathcal{L}$, as in Definition~\ref{d:approx-sys}.
      Let
      $\widehat{\Phi}: \widehat{\msc} \longrightarrow \widehat{\msc}$
      be a TPC endo-functor of systems and assume that
      $\widehat{\Phi}$ is
      $\{\mathcal{F}_1, \ldots, \mathcal{F}_l\}$-stabilizing. Then for
      every $\widetilde{L} \in \mathcal{L}$ we have
      $\widetilde{d}_{\text{int}}(\widetilde{L},
      \widetilde{\Phi}(\widetilde{L})) \leq \epsilon$.
  \end{prop}



